\hypertarget{sec:alcance}{}%
\begin{center}
    {\Large\bfseries 4. ALCANCE DE LA INVESTIGACIÓN}
\end{center}
\vspace{1em}

\noindent Este estudio es de tipo \textbf{descriptivo, con un componente comparativo}: busca caracterizar, a partir de una encuesta aplicada a los estudiantes de pregrado de la UPTC, Sede Tunja, el nivel de dependencia funcional hacia Nequi, DaviPlata y Bancolombia, las actividades cotidianas (alimentación, transporte y pagos académicos) que resultan afectadas cuando alguna de ellas falla y los mecanismos de contingencia que los estudiantes emplean, además de contrastar en qué medida esa dependencia varía entre las tres aplicaciones. Se elige este alcance, y no uno exploratorio, porque el fenómeno ya está definido y documentado: en 2023, según Colombia Fintech, 13 billeteras digitales reunieron 54,63 millones de usuarios en el país (cifra que incluye duplicidades, pues una misma persona usa varias billeteras) y procesaron 3.342 millones de transacciones por cerca de 332 billones de pesos, con Nequi como líder del segmento, seguida de Bancolombia a la Mano y DaviPlata~\cite{colfintech2024}, mientras que la caída total de los servicios digitales de Bancolombia y Nequi del 24 de octubre de 2025~\cite{elcolombiano2025} y la falla masiva de Bancolombia del 22 y 23 de febrero de 2026, que generó cerca de 8.000 quejas ante la Superfinanciera en 48 horas~\cite{infobae2026}, evidencian que las caídas del servicio son reales y recurrentes; tampoco se plantea un alcance explicativo, porque, como señala el numeral 1.1, la documentación local sobre el grado de dependencia estudiantil es escasa y el diseño no manipula ni controla las caídas del servicio, por lo que describe y compara, pero no atribuye causalidad.
\par\vspace{1.2em}

\noindent\textbf{Limitaciones de la investigación}
\par\vspace{0.6em}

\noindent Al basarse en una muestra no probabilística de estudiantes de una sola sede, los resultados describen a los participantes y no pueden generalizarse estadísticamente a toda la comunidad de la UPTC, a otras sedes ni a otras universidades; los datos provienen de autorreportes, por lo que la dependencia percibida puede diferir del uso real y los relatos sobre caídas pasadas están expuestos a sesgo de recuerdo; las interrupciones del servicio no se provocan ni se observan de forma controlada, y variables externas como la duración y el tipo de falla, la hora del día, el saldo disponible o la calidad de la conectividad no pueden controlarse del todo, de modo que el estudio no permite establecer relaciones causales; dado que un mismo estudiante puede usar más de una de las tres aplicaciones, las comparaciones entre ellas no se hacen entre grupos completamente independientes, lo que debe considerarse al elegir las pruebas estadísticas; y el instrumento de recolección se valida dentro de este mismo estudio, por lo que sus propiedades (validez de contenido y confiabilidad) deben reportarse junto con los resultados finales, no darse por supuestas.
\par\vspace{0.8em}